% Job posting: https://ca.linkedin.com/jobs/view/senior-software-engineer-ai-computer-vision-camera-systems-at-motorola-solutions-4417445051
\documentclass[11pt]{article}
\usepackage[left=1in,top=1in,right=1in,bottom=1in]{geometry}
\usepackage[hidelinks]{hyperref}
\usepackage{setspace}
\usepackage[T1]{fontenc}
\usepackage{newtxtext,newtxmath}
\ifdefined\pdfgentounicode
  \input{glyphtounicode}
  \pdfgentounicode=1
\fi
\setstretch{1.1}
\pagenumbering{gobble}
\begin{document}
\pagestyle{empty}

\noindent
Nishal Stanislaus Silva, PhD \\
Toronto, ON \\
nishal.silva@hotmail.com \,|\, +1 613 200 2030 \\[1em]

\noindent Dear Hiring Manager,

\vspace{0.5em}
\noindent I am writing to apply for the Senior Software Engineer -- AI/Computer Vision (Camera Systems) role at Motorola Solutions in Toronto. For 5.5 years at MAS Holdings, South Asia's largest apparel manufacturer, I designed and shipped camera-based computer-vision systems that ran continuously in production: an OpenCV template- and feature-matching pipeline for multilingual care-label QC, with an explainable defect overlay driving a PLC-controlled conveyor reject mechanism, cutting inspection time by 99.5\%, and a microscopic camera-array system for loom-side fabric defect detection that halved required operator headcount. During COVID-19, I also built and deployed a real-time camera-based remote-vitals monitoring device across hospital wards, recognized by Sri Lanka's Government Medical Officers' Association. This is the same discipline your camera systems team needs: turning CV models into hardware-integrated products that run reliably, in real time, at the point of capture.

\vspace{0.5em}
\noindent My PhD and postdoctoral work at the University of Trento sharpened that discipline under harder constraints. I owned the full ML pipeline for streaming sensor data end to end -- acquisition, feature engineering, training, and edge deployment -- and published a real-time detector running at 14 ms on a Raspberry Pi 4 (F1 0.76, 74x faster than a DTW baseline), backed by frozen-protocol benchmarks quantifying accuracy, latency, and CPU trade-offs. I work daily in C++17 for real-time engines and in Python (TensorFlow, PyTorch, OpenCV) for model development, and I am comfortable profiling and optimizing inference on ARM CPUs and embedded Linux -- exactly the toolchain camera-system analytics is built on.

\vspace{0.5em}
\noindent I have not worked in the security or surveillance industry specifically, and I want to be upfront about that. What transfers directly, though, is the core engineering discipline: shipping camera-based CV to production hardware under real latency and reliability constraints, and the rigor to benchmark and validate it before it ships. I am confident that discipline closes the gap quickly.

\vspace{0.5em}
\noindent I am a Canadian Permanent Resident based in Toronto, eligible to work without sponsorship, and available immediately. I would welcome the chance to discuss how my background in production camera-based CV and real-time edge inference can support Motorola Solutions' camera systems team.

\vspace{1em}
\noindent Sincerely, \\
Nishal Stanislaus Silva

\end{document}
